\documentclass[11pt]{article}

\usepackage[a4paper,margin=2.2cm]{geometry}
\usepackage[T1]{fontenc}
\usepackage[utf8]{inputenc}
\usepackage{lmodern}
\usepackage{amsmath,amssymb,mathtools}
\usepackage{microtype}
\usepackage{enumitem}
\usepackage{titlesec}
\usepackage[hidelinks]{hyperref}

\titleformat{\section}{\Large\bfseries}{\thesection}{0.75em}{}
\titleformat{\subsection}{\large\bfseries}{\thesubsection}{0.75em}{}
\titlespacing*{\section}{0pt}{1.4ex plus .2ex minus .2ex}{0.9ex}
\titlespacing*{\subsection}{0pt}{1.0ex plus .2ex minus .2ex}{0.5ex}

\newcommand{\R}{\mathbb{R}}
\newcommand{\diag}{\operatorname{diag}}
\newcommand{\rank}{\operatorname{rank}}

\begin{document}

\begin{center}
    {\LARGE\bfseries LFR Latent Dimension Reduction via SVD}\\[0.5em]
\end{center}

\vspace{0.7em}

\noindent
% The LPV-LFR realization derived in \texttt{LFR-derivation-supervisor.tex} uses
% $\Delta(Y) = Y I_6$, meaning six latent channels. This document derives a
% reduced realization with fewer channels, following the SVD-based procedure of
% R.~Tóth. Two methods are presented. Method~1 applies SVD directly to the full
% coupling matrix and is shown to give no reduction for the dual-gantry model.
% Method~2 first compresses the $D_{zw}$ channel and then applies the same SVD
% idea to the compressed loop.

\section{Method 1: SVD on the Full Coupling Matrix}
The stacked matrix is defined as
\[
S = \begin{bmatrix} B_w \\ I_6 \end{bmatrix} [C_z,\; D_{zw},\; D_{zu}] \in \mathbb{R}^{12\times 15}.
\]
The compact SVD of $S$ gives the minimal latent dimension $n = \operatorname{rank}(S)$.
\medskip
\\
\textbf{Step 1: Simplify $\operatorname{rank}(S)$}\\
The lower 6 rows of $S$ are exactly $[C_z,\; D_{zw},\; D_{zu}]$, so the row 
space of $[C_z,\; D_{zw},\; D_{zu}]$ is contained in the row space of $S$. 
Since $[C_z,\; D_{zw},\; D_{zu}]$ has only 6 rows, $\operatorname{rank}(S) \leq 6$.
Together:
\[
\operatorname{rank}(S) = \operatorname{rank}\bigl([C_z,\; D_{zw},\; D_{zu}]\bigr).
\]
\medskip
\\
\textbf{Step 2: Show $\operatorname{rank}([C_z,\; D_{zw},\; D_{zu}]) = 6$}\\
Extract the columns corresponding to the first block of $D_{zw}$ and $D_{zu}$,
forming the $6\times 6$ submatrix
\[
\begin{bmatrix} -M_0^{-1}M_1 & M_0^{-1} \\ I_3 & 0 \end{bmatrix}.
\]
Permuting the row blocks yields the block lower triangular matrix
\[
\begin{bmatrix} I_3 & 0 \\ -M_0^{-1}M_1 & M_0^{-1} \end{bmatrix},
\]
with $I_3$ and $M_0^{-1}$ on the diagonal. Since $M_0 \succ 0$ for all valid 
physical parameters, both blocks are invertible, so this submatrix is nonsingular.
\medskip
\\
\textbf{Conclusion}\\
A nonsingular $6\times 6$ submatrix exists $\Rightarrow$ 
$\operatorname{rank}(S) = 6$ $\Rightarrow$ Method 1 yields no reduction.

\newpage
\section{Method 2: SVD after Pre-Compression of $D_{zw}$}

\subsection{Setup}
The starting point is the derived LPV-LFR realization, with constant matrix $G$ and scheduling block $\Delta(Y) = Y I_6$. The state and loop equations are
\begin{align}
    \dot{x}(t) &= A_x\, x(t) + B_w\, w(t) + B_u\, u(t), 
    \label{eq:state2} \\
    z(t) &= C_z\, x(t) + D_{zw}\, w(t) + D_{zu}\, u(t),
    \label{eq:loopz2} \\
    w(t) &= Y\, I_6\, z(t),
    \label{eq:loopw2}
\end{align}
with
\begin{equation}
    C_z =
    \begin{bmatrix}
        -M_0^{-1}K & -M_0^{-1}C \\
        0 & 0
    \end{bmatrix},
    \quad
    D_{zw} =
    \begin{bmatrix}
        -M_0^{-1}M_1 & -M_0^{-1}M_2 \\
        I_3 & 0
    \end{bmatrix},
    \quad
    D_{zu} =
    \begin{bmatrix}
        M_0^{-1} \\
        0
    \end{bmatrix},
    \label{eq:method2CD}
\end{equation}
\begin{equation}
    B_w =
    \begin{bmatrix}
        0 & 0 \\
        -M_0^{-1}M_1 & -M_0^{-1}M_2
    \end{bmatrix},
    \quad
    B_u =
    \begin{bmatrix}
        0 \\ M_0^{-1}
    \end{bmatrix}.
    \label{eq:method2B}
\end{equation}
Method~1 gives no reduction because $[C_z,\; D_{zw},\; D_{zu}]$ has full 
row rank~$6$: $D_{zu} = [M_0^{-1};\, 0]$ introduces three directions 
independent of the bottom-row loop structure.
\\
Method~2 first compresses $D_{zw} = S_2 P$ via compact SVD, replacing $I_6$ 
by $P$ in the left factor of the stacked matrix:
\[
    S_{\mathrm{new}} =
    \begin{bmatrix} B_w \\[2pt] P \end{bmatrix}
    \begin{bmatrix} C_z & S_2 & D_{zu} \end{bmatrix}.
\]
The rank of the left factor is now determined by $\rank(P) = \rank(D_{zw})$.





\subsection{Structure and Rank of $D_{zw}$}
Using
\begin{equation}
    M_1 =
    \begin{bmatrix}
        0 & -m_h & 0 \\
        -m_h & 0 & 0 \\
        0 & 0 & 0
    \end{bmatrix},
    \qquad
    M_2 =
    \begin{bmatrix}
        0 & 0 & 0 \\
        0 & m_h & 0 \\
        0 & 0 & 0
    \end{bmatrix},
    \label{eq:M1M2}
\end{equation}
and writing the columns of $M_0^{-1}$ as
\begin{equation}
    M_0^{-1} = \begin{bmatrix} r_1 & r_2 & r_3 \end{bmatrix},
    \qquad
    \text{where } r_i \in \R^3 \text{ denotes the $i$th column of } M_0^{-1}.
    \label{eq:rcols}
\end{equation}
the matrix $D_{zw}$ from \eqref{eq:method2CD} decomposes into the two row blocks
\begin{equation}
    D_{zw} =
    \begin{bmatrix}
        D_{z_1} \\ D_{z_2}
    \end{bmatrix},
    \qquad
    D_{z_1} := \begin{bmatrix} -M_0^{-1}M_1 & -M_0^{-1}M_2 \end{bmatrix} \in \R^{3\times 6},
    \qquad
    D_{z_2} := \begin{bmatrix} I_3 & 0 \end{bmatrix} \in \R^{3\times 6}.
    \label{eq:DzwBlocks}
\end{equation}
Expanding $-M_0^{-1}M_1$ column by column,
\begin{equation}
    -M_0^{-1}M_1
    =
    \begin{bmatrix}
        -M_0^{-1}\!\begin{bmatrix} 0 \\ -m_h \\ 0 \end{bmatrix} &
        -M_0^{-1}\!\begin{bmatrix} -m_h \\ 0 \\ 0 \end{bmatrix} &
        -M_0^{-1}\!\begin{bmatrix} 0 \\ 0 \\ 0 \end{bmatrix}
    \end{bmatrix}
    =
    \begin{bmatrix}
        m_h r_2 & m_h r_1 & 0
    \end{bmatrix},
    \label{eq:M0invM1}
\end{equation}
and similarly
\begin{equation}
    -M_0^{-1}M_2
    =
    \begin{bmatrix}
        -M_0^{-1}\!\begin{bmatrix} 0 \\ 0 \\ 0 \end{bmatrix} &
        -M_0^{-1}\!\begin{bmatrix} 0 \\ m_h \\ 0 \end{bmatrix} &
        -M_0^{-1}\!\begin{bmatrix} 0 \\ 0 \\ 0 \end{bmatrix}
    \end{bmatrix}
    =
    \begin{bmatrix}
        0 & -m_h r_2 & 0
    \end{bmatrix}.
    \label{eq:M0invM2}
\end{equation}
Hence
\begin{equation}
    D_{z_1}
    =
    m_h
    \begin{bmatrix}
        r_2 & r_1 & 0 & 0 & -r_2 & 0
    \end{bmatrix}.
    \label{eq:Dz1structure}
\end{equation}
Let
\begin{equation}
    e_1 =
    \begin{bmatrix}
        1 \\ 0 \\ 0
    \end{bmatrix},
    \qquad
    e_2 =
    \begin{bmatrix}
        0 \\ 1 \\ 0
    \end{bmatrix},
    \qquad
    e_3 =
    \begin{bmatrix}
        0 \\ 0 \\ 1
    \end{bmatrix}
    \in \R^3
    \label{eq:basisvectors}
\end{equation}
denote the standard basis vectors of $\R^3$. Then the columns of $D_{zw}$ are
\begin{equation}
    d_1 =
    \begin{bmatrix}
        m_h r_2 \\ e_1
    \end{bmatrix},
    \qquad
    d_2 =
    \begin{bmatrix}
        m_h r_1 \\ e_2
    \end{bmatrix},
    \qquad
    d_3 =
    \begin{bmatrix}
        0 \\ e_3
    \end{bmatrix},
    \qquad
    d_4 =
    \begin{bmatrix}
        0 \\ 0
    \end{bmatrix},
    \qquad
    d_5 =
    \begin{bmatrix}
        -m_h r_2 \\ 0
    \end{bmatrix},
    \qquad
    d_6 =
    \begin{bmatrix}
        0 \\ 0
    \end{bmatrix}.
    \label{eq:Dzwcolumns}
\end{equation}
Therefore columns $4$ and $6$ are zero, so all nonzero columns of $D_{zw}$
lie among columns $1$, $2$, $3$, and $5$. Hence
\begin{equation}
    \rank(D_{zw}) \leq 4.
    \label{eq:rankDzwUpper}
\end{equation}
For the lower bound, note that columns 1, 2, 3, 4 and 5 are linearly independent:
\begin{equation}
    \alpha_1 d_1 + \alpha_2 d_2 + \alpha_3 d_3 + \alpha_5 d_5 = 0.
    \label{eq:lindepDzw}
\end{equation}
Looking at the lower three entries of \eqref{eq:lindepDzw} gives
\begin{equation}
    \alpha_1 e_1 + \alpha_2 e_2 + \alpha_3 e_3 = 0.
    \label{eq:lowerblock}
\end{equation}
Since $e_1,e_2,e_3$ are linearly independent, it follows that
\begin{equation}
    \alpha_1 = \alpha_2 = \alpha_3 = 0.
    \label{eq:alpha123zero}
\end{equation}
After \eqref{eq:alpha123zero}, the relation \eqref{eq:lindepDzw} reduces to
\begin{equation}
    \alpha_5 d_5 = 0.
\end{equation}
Using the expression for $d_5$ from \eqref{eq:Dzwcolumns}, this becomes
\begin{equation}
    \alpha_5
    \begin{bmatrix}
        -m_h r_2 \\ 0
    \end{bmatrix}
    = 0.
    \label{eq:alpha5eq} 
\end{equation}
Since $m_h > 0$ and $r_2$ is the second column of the invertible matrix $M_0^{-1}$, we have $r_2 \neq 0$. Therefore \eqref{eq:alpha5eq} implies
\begin{equation}
    \alpha_5 = 0.
\end{equation}
Hence columns $1$, $2$, $3$, and $5$ are linearly independent, so
\begin{equation}
    \rank(D_{zw}) \geq 4.
    \label{eq:rankDzwLower}
\end{equation}
Combining \eqref{eq:rankDzwUpper} and \eqref{eq:rankDzwLower} yields
\begin{equation}
    \rank(D_{zw}) = 4.
    \label{eq:rankDzw}
\end{equation}





\subsection{Structural Identity $B_w = E D_{zw}$}
Comparing \eqref{eq:method2B} and \eqref{eq:DzwBlocks}, $B_w$ places $D_{z_1}$ in its lower block and zeros in its upper block. Define
\begin{equation}
    E :=
    \begin{bmatrix}
        0_{3\times 3} & 0_{3\times 3} \\
        I_3 & 0_{3\times 3}
    \end{bmatrix}
    \in \R^{6\times 6}.
    \label{eq:Edef}
\end{equation}
Then
\begin{equation}
    E D_{zw}
    =
    \begin{bmatrix}
        0 \\
        D_{z_1}
    \end{bmatrix}
    =
    B_w,
    \label{eq:BwEDzw}
\end{equation}
so
\begin{equation}
    B_w = E D_{zw}.
\end{equation}

\subsection{compact SVD of $D_{zw}$ and Loop Pre-Compression}
Since $\rank(D_{zw}) = 4$, compute the compact SVD
\begin{equation}
    D_{zw} = \hat{U}\,\hat{\Sigma}\,\hat{V}^\top,
    \qquad
    \hat{U} \in \R^{6\times 4},\quad
    \hat{\Sigma} \in \R^{4\times 4},\quad
    \hat{V}^\top \in \R^{4\times 6},
    \label{eq:DzwSVD}
\end{equation}
and define
\begin{equation}
    S_2 := \hat{U}\,\hat{\Sigma} \in \R^{6\times 4},
    \qquad
    P := \hat{V}^\top \in \R^{4\times 6}, \qquad \rank(P) = 4,
    \label{eq:S2P}
\end{equation}
so that
\begin{equation}
    D_{zw} = S_2\,P.
    \label{eq:Dzwfactor}
\end{equation}
$D_{zw}w$ and $B_w w$ both depend on $w$ only through $Pw$. From \eqref{eq:BwEDzw},
\begin{equation}
    B_w = E D_{zw} = E S_2 P,
    \label{eq:BwES2P_repeat}
\end{equation}
motivating the compressed loop variable
\begin{equation}
    \hat{w}(t) := P\,w(t) \in \R^4.
    \label{eq:whatdef}
\end{equation}
To obtain a closed loop with the same scalar scheduling structure, define
\begin{equation}
    \hat{z}(t) := P\,z(t) \in \R^4.
    \label{eq:zhatdef}
\end{equation}
Then the original scheduling relation $w(t)=Y I_6 z(t)$ projects to
\begin{equation}
    \hat{w}(t)
    = P w(t)
    = P\bigl(Y I_6 z(t)\bigr)
    = Y\,P z(t)
    = Y\,\hat{z}(t)
    = Y I_4 \hat{z}(t),
    \label{eq:whatnew}
\end{equation}
Hence the projected loop
retains the same repeated-scalar structure, but now with four channels instead
of six.
\\
Therefore the realization can be rewritten entirely in terms of the compressed loop variable $\hat{w}(t)$:
\begin{equation}
    D_{zw} w(t) = S_2 \hat{w}(t),
    \qquad
    B_w w(t) = E S_2 \hat{w}(t).
    \label{eq:compressedwterms}
\end{equation}
\subsection{Modified Stacked Matrix}

Substituting \eqref{eq:compressedwterms} into \eqref{eq:loopz2} gives
\begin{equation}
    z(t) = C_z\, x(t) + S_2\, \hat{w}(t) + D_{zu}\, u(t).
    \label{eq:loopzcompressed}
\end{equation}
Multiplying \eqref{eq:loopzcompressed} by $P$ yields
\begin{equation}
    \hat{z}(t) = P C_z\, x(t) + P S_2\, \hat{w}(t) + P D_{zu}\, u(t).
    \label{eq:zhatloop}
\end{equation}
Using \eqref{eq:whatnew} and \eqref{eq:loopzcompressed},
\begin{equation}
    \hat{w}(t)
    = Y\, P\,[C_z,\; S_2,\; D_{zu}]
      \begin{bmatrix}
          x(t) \\ \hat{w}(t) \\ u(t)
      \end{bmatrix},
    \label{eq:whatstack}
\end{equation}
and the state equation \eqref{eq:state2} gives
\begin{equation}
    B_w\, w(t)
    = Y\, B_w\,[C_z,\; S_2,\; D_{zu}]
      \begin{bmatrix}
          x(t) \\ \hat{w}(t) \\ u(t)
      \end{bmatrix}.
    \label{eq:Bwwstack}
\end{equation}
Stacking \eqref{eq:whatstack} and \eqref{eq:Bwwstack} yields
\begin{equation}
    Y \underbrace{
        \begin{bmatrix}
            B_w \\[2pt] P
        \end{bmatrix}
        [C_z,\; S_2,\; D_{zu}]
    }_{S_{\mathrm{new}}\;\in\;\R^{10\times 13}}
    \begin{bmatrix}
        x(t) \\ \hat{w}(t) \\ u(t)
    \end{bmatrix}
    =
    \begin{bmatrix}
        B_w\, w(t) \\ \hat{w}(t)
    \end{bmatrix},
    \label{eq:Snew}
\end{equation}

\subsection{Rank of the Compressed Stacked Matrix}
From \eqref{eq:BwES2P_repeat}, $B_w = ES_2P$, so every row of $B_w$ is a 
linear combination of rows of $P$, and therefore $\operatorname{row}(B_w) 
\subseteq \operatorname{row}(P)$. Stacking adds no new row directions, so
\begin{equation}
    \rank\!\begin{bmatrix} B_w \\ P \end{bmatrix} = \rank(P) = 4.
    \label{eq:rankBwP}
\end{equation}
To determine rank$(S_\text{new})$ exactly, we show that the right factor
\begin{equation}
    F := [C_z,\; S_2,\; D_{zu}] \in \R^{6\times 13}
    \label{eq:Fdef}
\end{equation}
has full row rank. Let
\begin{equation}
    J := \begin{bmatrix} 0 & I_3 \end{bmatrix} \in \R^{3\times 6}
    \label{eq:Jdef}
\end{equation}
select the lower three rows of any $6\times n$ matrix. From 
\eqref{eq:DzwBlocks}, $J D_{zw} = D_{z_2} = [I_3\ \ 0]$, which has rank~$3$. 
Substituting $D_{zw} = S_2 P$ and applying the rank product inequality,
\begin{equation}
    3 = \rank(J D_{zw}) = \rank\!\bigl((J S_2) P\bigr) \leq \rank(J S_2) \leq 3,
    \label{eq:rankJS2}
\end{equation}
so $\rank(J S_2) = 3$, meaning three columns of $S_2$ have linearly independent 
lower three entries. Denote by $A \in \R^{3\times 3}$ the lower three entries 
of those columns (invertible by construction) and by $* \in \R^{3\times 3}$ 
the upper three entries. Together with the three columns of $D_{zu}$, the 
corresponding $6\times 6$ submatrix of $F$ is
\begin{equation}
    \begin{bmatrix} * & M_0^{-1} \\ A & 0 \end{bmatrix}.
    \label{eq:submat}
\end{equation}
Permuting the row blocks yields the block lower triangular matrix
\begin{equation}
    \begin{bmatrix} A & 0 \\ * & M_0^{-1} \end{bmatrix},
    \label{eq:submatperm}
\end{equation}
with $A$ and $M_0^{-1}$ on the diagonal. Both are invertible: $A$ by 
construction and $M_0^{-1}$ since $M_0 \succ 0$ for all valid physical 
parameters. Therefore this submatrix is nonsingular, which gives



\begin{equation}
    \rank(F) = 6.
    \label{eq:rankF}
\end{equation}
It follows that $\rank(S_{\mathrm{new}}) = 4$ by the following two-sided
argument. For the upper bound, the rank product inequality gives
\begin{equation}
    \rank(S_{\mathrm{new}})
    = \rank\!\left(
        \begin{bmatrix} B_w \\ P \end{bmatrix} F
      \right)
    \leq
    \rank\!\begin{bmatrix} B_w \\ P \end{bmatrix}
    = 4.
    \label{eq:rankSnewUpper}
\end{equation}
For the lower bound, since $F$ has full row rank there exists a right inverse
$F^\dagger$ satisfying $F F^\dagger = I_6$. Therefore
\begin{equation}
    \begin{bmatrix} B_w \\ P \end{bmatrix}
    = S_{\mathrm{new}}\, F^\dagger,
    \label{eq:rightinverse}
\end{equation}
and applying the rank product inequality again,
\begin{equation}
    4 = \rank\!\begin{bmatrix} B_w \\ P \end{bmatrix}
    \leq \rank(S_{\mathrm{new}}).
    \label{eq:rankSnewLower}
\end{equation}
Combining \eqref{eq:rankSnewUpper} and \eqref{eq:rankSnewLower},
\begin{equation}
    \rank(S_{\mathrm{new}}) = 4.
    \label{eq:rankSnew}
\end{equation}


\subsection{Compact SVD of $S_{\mathrm{new}}$ and Reduced Matrices}

Since $\rank(S_{\mathrm{new}}) = 4$, compute the compact SVD
\begin{equation}
    S_{\mathrm{new}} = U_{\mathrm{new}}\,\Sigma_{\mathrm{new}}\,V_{\mathrm{new}}^\top,
    \qquad
    \Sigma_{\mathrm{new}} = \diag(\sigma_1, \sigma_2, \sigma_3, \sigma_4),
    \qquad
    \sigma_i > 0.
    \label{eq:SnewSVD}
\end{equation}
Define
\begin{equation}
    L_{\mathrm{new}} := U_{\mathrm{new}}\,\Sigma_{\mathrm{new}}^{1/2},
    \qquad
    R_{\mathrm{new}} := \Sigma_{\mathrm{new}}^{1/2}\,V_{\mathrm{new}}^\top,
    \label{eq:LRnew}
\end{equation}
so that $S_{\mathrm{new}} = L_{\mathrm{new}} R_{\mathrm{new}}$ by construction.
Since $S_{\mathrm{new}}$ has 10 rows corresponding to the stacking of $B_w$
(6 rows) and $P$ (4 rows), split $L_{\mathrm{new}}$ as
\begin{equation}
    L_{\mathrm{new}} =
    \begin{bmatrix} L_x \\ L_z \end{bmatrix},
    \qquad L_x \in \R^{6\times 4},\quad L_z \in \R^{4\times 4},
    \label{eq:Lsplit}
\end{equation}
and since $R_{\mathrm{new}}$ has 13 columns corresponding to $x$ (6), $\hat{w}$
(4), and $u$ (3), split $R_{\mathrm{new}}$ as
\begin{equation}
    R_{\mathrm{new}} = \bigl[R_x,\; R_{\hat{w}},\; R_u\bigr],
    \qquad R_x \in \R^{4\times 6},\quad R_{\hat{w}} \in \R^{4\times 4},
    \quad R_u \in \R^{4\times 3}.
    \label{eq:Rsplit}
\end{equation}
Define the reduced latent variable
\begin{equation}
    \tilde{z}(t)
    := R_{\mathrm{new}}
    \begin{bmatrix}
        x(t) \\ \hat{w}(t) \\ u(t)
    \end{bmatrix}
    \in \R^4.
    \label{eq:ztilde2}
\end{equation}
Substituting $\tilde{w}(t) = YI_4\tilde{z}(t)$ and $B_w w(t) = L_x \tilde{w}(t)$
into the state and loop equations, the reduced constant matrix $\tilde{G}$ has 
entries
\begin{equation}
    \tilde{B}_w = L_x,\quad
    \tilde{C}_z = R_x,\quad
    \tilde{D}_{zw} = R_{\hat{w}}\, L_z,\quad
    \tilde{D}_{zu} = R_u,
    \label{eq:Gtilde2}
\end{equation}
with $A_x$, $B_u$, $C_y$, $D_{yw}$, $D_{yu}$ unchanged, and the scheduling
block reduced to $\Delta(Y) = Y I_4$.
\\
The numerical values of $\tilde{G}$ depend on the SVDs in \eqref{eq:DzwSVD}
and \eqref{eq:SnewSVD}. However, the reduction from six to four latent channels
is structural: it follows from $\rank(D_{zw}) = 4$, which holds for all
physically valid parameter values ($m_h > 0$, $M_0 \succ 0$).
%kijken matlab toepassing laten zien!!!

%%%%%%%%%%%%%%%%%%%%%%%%%%%%%%%%%%%%%%%%%%%%%%%%%%%
% Moet nog gedubbel checkt worden
\newpage
\section{Preservation of well-posedness under latent dimension reduction}

We show that the reduced latent realization is well posed for exactly the same
values of $Y$ as the original realization. Hence well-posedness only needs to
be proved once for the original loop matrix.

\paragraph{Original and pre-compressed loop matrices.}
Since $\Delta(Y)=YI_6$, the original loop equation is governed by
\begin{equation}
    L_6(Y) := I_6 - YD_{zw}.
    \label{eq:L6_def}
\end{equation}
Using the factorization $D_{zw}=S_2P$, this becomes
\begin{equation}
    L_6(Y)=I_6-YS_2P.
    \label{eq:L6_factorized}
\end{equation}
From \eqref{eq:zhatloop} and \eqref{eq:whatnew}, the pre-compressed loop is
\begin{equation}
    \hat L_4(Y) := I_4 - YPS_2.
    \label{eq:L4hat_def}
\end{equation}

\paragraph{Equivalence of the original and pre-compressed loop.}
Using the identity $\det(I_m+AB)=\det(I_n+BA)$ for compatible rectangular
matrices, with
\[
A=-YS_2 \in \R^{6\times 4},
\qquad
B=P \in \R^{4\times 6},
\]
we obtain
\begin{equation}
    \det(I_6-YS_2P)=\det(I_4-YPS_2).
    \label{eq:WA_identity}
\end{equation}
Hence
\begin{equation}
    L_6(Y)\ \text{is invertible}
    \iff
    \hat L_4(Y)\ \text{is invertible}.
    \label{eq:L6_L4hat_equiv}
\end{equation}

\paragraph{Final reduced loop matrix.}
The reduced realization uses
\begin{equation}
    \tilde D_{zw} = R_{\hat w}L_z,
    \qquad
    \tilde L_4(Y) := I_4 - Y\tilde D_{zw}.
    \label{eq:L4tilde_def}
\end{equation}
From the lower block of the factorization
\[
S_{\mathrm{new}} = L_{\mathrm{new}}R_{\mathrm{new}},
\qquad
L_{\mathrm{new}}=
\begin{bmatrix}
L_x\\
L_z
\end{bmatrix},
\qquad
R_{\mathrm{new}}=[R_x,\;R_{\hat w},\;R_u],
\]
and the definition of $S_{\mathrm{new}}$, we have
\begin{equation}
    L_z[R_x,\;R_{\hat w},\;R_u]
    =
    [PC_z,\;PS_2,\;PD_{zu}],
    \label{eq:lower_block_identity}
\end{equation}
so in particular
\begin{equation}
    L_zR_{\hat w}=PS_2.
    \label{eq:LzRhatw_identity}
\end{equation}

\paragraph{Invertibility of $L_z$.}
Let
\[
F := [C_z,\;S_2,\;D_{zu}] \in \R^{6\times 13}.
\]
Since $\rank(F)=6$, there exists a right inverse $F^\dagger$ such that
$FF^\dagger=I_6$. The lower four rows of $S_{\mathrm{new}}$ are
\[
PF := [PC_z,\;PS_2,\;PD_{zu}].
\]
Because
\[
P=(PF)F^\dagger,
\]
it follows that
\[
\rank(P)\le \rank(PF).
\]
Since also $\rank(PF)\le \rank(P)=4$, we obtain
\begin{equation}
    \rank(PF)=4.
    \label{eq:rankPF}
\end{equation}
Using \eqref{eq:lower_block_identity},
\[
PF=L_zR_{\mathrm{new}}.
\]
Now $R_{\mathrm{new}}\in\R^{4\times 13}$ has full row rank $4$, because
$S_{\mathrm{new}}=L_{\mathrm{new}}R_{\mathrm{new}}$ is a rank-$4$ factorization.
Hence \eqref{eq:rankPF} implies
\[
\rank(L_z)=4.
\]
Therefore
\begin{equation}
    L_z \in \R^{4\times 4}
    \quad \text{is invertible.}
    \label{eq:Lz_invertible}
\end{equation}

\paragraph{Similarity of the reduced and pre-compressed loop.}
From \eqref{eq:LzRhatw_identity} and \eqref{eq:Lz_invertible},
\[
R_{\hat w}=L_z^{-1}PS_2.
\]
Substituting this into \eqref{eq:L4tilde_def} gives
\begin{equation}
    \tilde D_{zw}
    = R_{\hat w}L_z
    = L_z^{-1}(PS_2)L_z.
    \label{eq:Dzwtilde_similarity}
\end{equation}
Therefore
\begin{equation}
    \tilde L_4(Y)
    = I_4 - Y\tilde D_{zw}
    = L_z^{-1}(I_4-YPS_2)L_z,
    \label{eq:L4tilde_similarity}
\end{equation}
and hence
\begin{equation}
    \tilde L_4(Y)\ \text{is invertible}
    \iff
    \hat L_4(Y)\ \text{is invertible}.
    \label{eq:L4tilde_L4hat_equiv}
\end{equation}

\paragraph{Conclusion.}
Combining \eqref{eq:L6_L4hat_equiv} and \eqref{eq:L4tilde_L4hat_equiv}, we obtain
\begin{equation}
    I_6-YD_{zw}\ \text{is invertible}
    \iff
    I_4-Y\tilde D_{zw}\ \text{is invertible}.
\end{equation}
Thus the final reduced latent realization is well posed for exactly the same
values of $Y$ as the original realization. Consequently, well-posedness only
needs to be proved once for the original loop matrix.

\end{document}
https://chatgpt.com/c/69ddf365-c818-8391-bf92-7cd24efe21f7
Using the Weinstein--Aronszajn identity $\det(I_m+AB)=\det(I_n+BA)$ for compatible rectangular matrices~\cite{waterloo_linear_algebra_notes}, ...

@misc{waterloo_linear_algebra_notes,
  author       = {Chi, Lap},
  title        = {Linear Algebra Background},
  institution  = {University of Waterloo},
  note         = {Fact 2.28 (Weinstein--Aronszajn Identity)},
  year         = {n.d.},
  url          = {https://cs.uwaterloo.ca/~lapchi/cs466/notes/linear-algebra-background.pdf}
}